L'annexe ci-dessous contient toutes les fonctions demandées dans les questions qui n'apparaissent pas explicitement dans le chapitre "Réponses aux questions théoriques" (cf. remarques en début de chaque partie du chapitre sus-mentionné). Nous avons pris la liberté d'ajouter quelques fonctions supplémentaires pour en alléger d'autres.

\begin{codeblock}[fonctions.hpp]
#ifndef FONCTIONS_H
#define FONCTIONS_H

#include <tuple>
#include <vector>


#include "triangle.hpp"

using namespace std;

typedef vector<vector<double>> matrix;
typedef tuple<vector<double>,vector<double>> duplix;
typedef vector<Triangle> maillage;

vector<double> Subdiv(double a, int N);

int numgb(int N, int M, int i, int j);

tuple<int,int> invnumgb(int N, int M,int s);

maillage maillageTR(int N, int M);

duplix CoordsTrig(double a, double b, int N,int M, Triangle T);

matrix CalcMatBT(vector<double> xs, vector<double> ys);

vector<double> integ_eta_triang(double (* eta)(double,double),
    maillage TRG,int N, int M,double a,double b);

matrix DiffTerm(duplix xs_ys, double val);

matrix ReacTerm(duplix xs_ys);

vector<double> matvec(vector<double> V, maillage TRG, int N, int M,
    double a, double b, double (* eta)(double,double));

#endif

vector<double> scdmembre(double rhsf(double,double, int), int N, int M, maillage TRG, double a, double b, int m);

double normL2(vector<double> V,maillage TRG, int N, int M, double a, double b);

double normL2Grad(vector<double> V, maillage TRG, int N, int M, double a, double b);

double pdt_sc(vector<double> u, vector<double> v);

vector<double> cl_vec(double lbd, vector<double> u, double mu, vector<double> v);

double max_abs(vector<double> V);

vector<double> bicg_stab(vector<double> B,maillage TRG, int N, int M, double a,
    double b, double tol,int max_it, double (* eta) (double,double));

vector<double> erreurs(double (* solExa) (double,double, int),vector<double> uh,maillage TRG, int N, int M, double a, double b, int m);
\end{codeblock}

\begin{codeblock}[fonctions.cpp]
#include <stdlib.h>
#include <cmath>

#include "fonctions.hpp"



// Fonction qui renvoie une subdivision uniforme de [-a,a]
vector<double> Subdiv(double a, int N){
    vector<double> subdiv = vector<double>(N+1);
    double pas = 2*a / N;

    for (int k = 0; k <= N; k++) subdiv[k] = -a + k*pas;

    return subdiv;
}

// Fonction qui renvoie le numéro associé à des coordonnées
int numgb(int N, int M, int i, int j){ return (N+1)*j + i; }

//Fonction qui étant donné entier s valide, donne les coordonnées du noeud
tuple<int,int> invnumgb(int N, int M,int s){
    return {s % (N+1),s / (N+1)};
}

//Fonction qui étant donné un pavé de NM rectangles renvoie un TRG
//triangulaire sous forme d'une liste qui contient les triangles
maillage maillageTR(int N, int M){
    maillage TRG(2*N*M);
    int k = 0;
    bool direction = true; // true pour "/" et false pour "\"

    for (int l = 0; l < M; l++){
        for (int c = 0; c < N ; c++){

            //Pour alterner le sens des triangles à chaque petit rectangle
            direction = (l+c)%2 == 0;

            // Calcul des coordonnées des quatres sommets du rectangle
            int bas_g = numgb(N,M,c,l);
            int bas_d = numgb(N,M,c+1,l);
            int haut_g = numgb(N,M,c,l+1);
            int haut_d = numgb(N,M,c+1,l+1);

            // Création des triangles
            if (direction){
                TRG[k] = Triangle(bas_g,haut_g,haut_d);
                TRG[k+1] = Triangle(bas_g,bas_d,haut_d);
            } else {
                TRG[k] = Triangle(bas_g,haut_g,bas_d);
                TRG[k+1] = Triangle(haut_g,haut_d,bas_d);
            }

            k += 2;
        }
    }

    return TRG;
}


// Fonction qui renvoie les coordonnées des points du triangle
duplix CoordsTrig(double a, double b, int N,
    int M, Triangle T){
    vector<double> xs = vector<double>(3);
    vector<double> ys = vector<double>(3);

    for (int k = 0; k < 3; k++){
        tuple<int,int> tup = invnumgb(N,M,T.get(k));
        xs[k] = -a + (2*a)/N * get<0>(tup);
        ys[k] = -b + (2*b)/M * get<1>(tup);

    }

    return {xs,ys};
}


// Fonction qui renvoie la jacobienne de F_T
// Voir question 20, elle justifie ce calcul
matrix CalcMatBT(vector<double> xs, vector<double> ys){
    double bt_00 = xs[1] - xs[0];
    double bt_10 = ys[1] - ys[0];
    double bt_01 = xs[2] - xs[0];
    double bt_11 = ys[2] - ys[0];

    return {{bt_00,bt_01},{bt_10,bt_11}};
}




double integ_triangle(double (* f)(double,double), Triangle T,int N, int M,double a,double b){

    // Calcul du determinant de B_T
    duplix xs_ys = CoordsTrig(a,b,N,M,T);
    vector<double> xs = get<0>(xs_ys);
    vector<double> ys = get<1>(xs_ys);

    double det = abs((xs[1]-xs[0])*(ys[2]-ys[0]) - (ys[1]-ys[0])*(xs[2]-xs[0]));

    // On determine les milieux des côtes de T
    double m0x = (xs[0] + xs[1])/2.;
    double m0y = (ys[0] + ys[1])/2.;
    double m1x = (xs[1] + xs[2])/2.;
    double m1y = (ys[1] + ys[2])/2.;
    double m2x = (xs[0] + xs[2])/2.;
    double m2y = (ys[0] + ys[2])/2.;

    // Calculs de f \circ F_T en les milieux des côtés du triangles T^
    double f0 = f(m0x,m0y);
    double f1 = f(m1x,m1y);
    double f2 = f(m2x,m2y);

    return 1./6. * det * (f0 + f1 + f2);
}


vector<double> integ_eta_triang(double (* eta)(double,double),
    maillage TRG,int N, int M,double a,double b){

    vector<double> ET = vector<double>(TRG.size());
    int k =0;

    // Parcours du maillage
    for (Triangle T : TRG){


        ET[k] = integ_triangle(eta,T,N,M,a,b);
        k++;
    }

    return ET;
}


matrix DiffTerm(duplix xs_ys, double val){
    vector<double> xs = get<0>(xs_ys);
    vector<double> ys = get<1>(xs_ys);

    matrix BT = CalcMatBT(xs,ys);
    double det = (xs[1]-xs[0])*(ys[2]-ys[0]) - (ys[1]-ys[0])*(xs[2]-xs[0]);

    //Construction de la transposée de l'inverse de BT
    matrix trsp_inv_BT = {{0,0},{0,0}};
    trsp_inv_BT[0][0] = 1/det * BT[1][1];
    trsp_inv_BT[0][1] = -1/det * BT[1][0];
    trsp_inv_BT[1][0] = -1/det * BT[0][1];
    trsp_inv_BT[1][1] = 1/det * BT[0][0];


    //Calcul des Nabla_lambda^
    vector<double> NLcha0 = {{-1,-1}}; //lambdacha0 = (1-x^-y^)
    vector<double> NLcha1 = {{1,0}}; //lambdacha1 = x^
    vector<double> NLcha2 = {{0,1}}; //lambdacha2 = y^
    //Calcul des Nabla_lambda = BT*Nabla_lambdacha
    double nl00 = trsp_inv_BT[0][0] * NLcha0[0] + trsp_inv_BT[0][1] * NLcha0[1];
    double nl01 = trsp_inv_BT[1][0] * NLcha0[0] + trsp_inv_BT[1][1] * NLcha0[1];

    double nl10 = trsp_inv_BT[0][0] * NLcha1[0] + trsp_inv_BT[0][1] * NLcha1[1];
    double nl11 = trsp_inv_BT[1][0] * NLcha1[0] + trsp_inv_BT[1][1] * NLcha1[1];

    double nl20 = trsp_inv_BT[0][0] * NLcha2[0] + trsp_inv_BT[0][1] * NLcha2[1];
    double nl21 = trsp_inv_BT[1][0] * NLcha2[0] + trsp_inv_BT[1][1] * NLcha2[1];

    //Calcul des produits scalaires * ET
    double val00 = val * (nl00 * nl00 + nl01 * nl01);
    double val01 = val * (nl00 * nl10 + nl01 * nl11);
    double val02 = val * (nl00 * nl20 + nl01 * nl21);
    double val11 = val * (nl10 * nl10 + nl11 * nl11);
    double val12 = val * (nl10 * nl20 + nl11 * nl21);
    double val22 = val * (nl20 * nl20 + nl21 * nl21);

    return {{val00, val01, val02},
        {val01, val11, val12},
        {val02, val12, val22}
    };
}


matrix ReacTerm(duplix xs_ys){
    vector<double> xs = get<0>(xs_ys);
    vector<double> ys = get<1>(xs_ys);

    double det = (xs[1]-xs[0])*(ys[2]-ys[0]) - (ys[1]-ys[0])*(xs[2]-xs[0]);

    //Calcul des integrales des lambda_i * lambda_j
    double int_lambda_ii = abs(det) * 1.0/12.0;
    double int_lambda_ij = abs(det) * 1.0/24.0;

    return {{int_lambda_ii, int_lambda_ij, int_lambda_ij},
        {int_lambda_ij, int_lambda_ii, int_lambda_ij},
        {int_lambda_ij, int_lambda_ij,int_lambda_ii}};
}




vector<double> matvec(vector<double> V, maillage TRG, int N, int M,
    double a, double b, double (* eta)(double,double)){

    int K = 2*N*M;
    vector<double> W((N+1)*(M+1));

    // Contient les valeurs de l'integrales de eta sur chaque T
    vector<double> VALS = integ_eta_triang(eta,TRG, N,M,a,b);

    for (int t = 0; t < K; t++){
        Triangle T = TRG[t];

        // Calcul de B_T
        duplix xs_ys = CoordsTrig(a,b,N,M,T);
        matrix D = DiffTerm(xs_ys,VALS[t]);
        matrix R = ReacTerm(xs_ys);

        for (int i = 0; i < 3; i++){
            int k = T.get(i);

            // Calcul de la contribution
            double res = 0;

            for (int r = 0; r < 3; r++){
                int j = T.get(r);

                // aT = Diff + Reac
                double PROD2 = D[r][i]+R[r][i];
                res += V[j] * PROD2;
            }

            W[k] += res;
        }
    }

    return W;
}



vector<double> scdmembre(double rhsf(double,double,int), int N, int M, maillage TRG,
    double a, double b, int m){

    vector<double> B((N+1)*(M+1));
    int K = 2*N*M;

    for (int t = 0; t < K; t++){

        duplix xs_ys = CoordsTrig(a,b,N,M,TRG[t]);
        vector<double> xs = get<0>(xs_ys);
        vector<double> ys = get<1>(xs_ys);

        double det = abs((xs[1]-xs[0])*(ys[2]-ys[0]) - (ys[1]-ys[0])*(xs[2]-xs[0]));

        // calcul des milieux
        double m0x = (xs[0] + xs[1])/2.;
        double m0y = (ys[0] + ys[1])/2.;
        double m1x = (xs[1] + xs[2])/2.;
        double m1y = (ys[1] + ys[2])/2.;
        double m2x = (xs[0] + xs[2])/2.;
        double m2y = (ys[0] + ys[2])/2.;

        // Calculs de (f \circ F_T)*(w \circ F_T) en les milieux des côtés du triangles T^
        double f0 = rhsf(m0x,m0y,m);
        double f1 = rhsf(m1x,m1y,m);
        double f2 = rhsf(m2x,m2y,m);

        // Calcul des contributions des sommets de T aux noeuds
        B[TRG[t].get(0)] += (det * (f0*0.5 + f2*0.5))/6;
        B[TRG[t].get(1)] += (det * (f0*0.5 + f1*0.5))/6;
        B[TRG[t].get(2)] += (det * (f1*0.5 + f2*0.5))/6;
    }

    return B;
}


double normL2(vector<double> V,maillage TRG, int N, int M, double a, double b){

    int K = 2*N*M;
    int G = (N+1)*(M+1);
    vector<double> W(G);

    // Calcul de AV
    for (int t = 0; t < K; t++){
        Triangle T = TRG[t];

        duplix xs_ys = CoordsTrig(a,b,N,M,T);
        matrix R = ReacTerm(xs_ys);

        for (int i = 0; i < 3; i++){
            int k = T.get(i);

            // Calcul de la contribution
            double res = 0;

            for (int r = 0; r < 3; r++){
                int j = T.get(r);

                // aT = Diff + Reac (mais Diff = 0 car eta = 0)
                res += V[j] * R[r][i];
            }

            W[k] += res;
        }
    }

    // Calcul de V^T * (AV)
    double norme_carre = 0;
    for (int k = 0; k < G; k++)
        norme_carre += V[k]*W[k];

    return sqrt(norme_carre);
}


double normL2Grad(vector<double> V, maillage TRG, int N, int M, double a,
    double b){

    int K = 2*N*M;
    int G = (N+1)*(M+1);
    vector<double> W(G);

    // Contient les valeurs de l'integrales de eta sur chaque T
    vector<double> VALS = integ_eta_triang([](double,double){return 1.;},TRG, N,M,a,b);

    // Calcul de AV
    for (int t = 0; t < K; t++){
        Triangle T = TRG[t];

        duplix xs_ys = CoordsTrig(a,b,N,M,T);
        matrix D = DiffTerm(xs_ys,VALS[t]);

        for (int i = 0; i < 3; i++){
            int k = T.get(i);

            // Calcul de la contribution
            double res = 0;

            for (int r = 0; r < 3; r++){
                int j = T.get(r);

                // aT = Diff + Reac (mais reac = 0 car lambda = 0)
                res += V[j] * D[r][i];
            }

            W[k] += res;
        }
    }

    // Calcul de V^T * (AV)
    double norme_carre = 0;
    for (int k = 0; k < G; k++)
        norme_carre += V[k]*W[k];

    return sqrt(norme_carre);
}

// Fonction qui calcule le produit scalaire de deux vecteurs
double pdt_sc(vector<double> u, vector<double> v){
    double prod = 0;
    int N = u.size();

    for (int k = 0; k < N; k++){
        prod += u[k]*v[k];
    }

    return prod;
}

// Fonction qui renvoie la combinaison linéaire de deux vecteurs
vector<double> cl_vec(double lbd, vector<double> u, double mu, vector<double> v){
    int L = u.size();
    vector<double> cl(L);

    for (int k = 0; k < L; k++){
        cl[k] = lbd*u[k] + mu*v[k];
    }

    return cl;
}


// Fonction qui renvoie la composante max en valeur absolue d'un vecteur
double max_abs(vector<double> V){
    double max = V[0];

    for (double x : V)
        if (abs(x) > max) max = abs(x);

    return max;
}

vector<double> bicg_stab(vector<double> B,maillage TRG, int N, int M, double a,
    double b, double tol,int max_it, double (* eta) (double,double)){

    vector<double> AW,AS,S,R_tmp;
    vector<double> X((N+1)*(M+1),0);

    // on choisit X0=0 (ca nous évite un prod mat vec)
    vector<double> R = B;
    vector<double> R_etoile = R;


    vector<double> W = R;
    int it = 0;

    // Itérations, bicg_stab selon l'algorithme donné
    while (max_abs(R) > tol && it < max_it){
        AW = matvec(W,TRG,N,M,a,b,eta);
        double alpha = pdt_sc(R,R_etoile) / pdt_sc(AW,R_etoile);

        S = cl_vec(1.,R,-alpha,AW);

        AS = matvec(S,TRG,N,M,a,b,eta);
        double omega = pdt_sc(AS,S)/pdt_sc(AS,AS);

        // Actualisation des variables pour le prochain tour de boucle
        X = cl_vec(1.,X,1.,cl_vec(alpha,W,omega,S));
        R_tmp = R;
        R = cl_vec(1.,S,-omega,AS);

        double beta = (pdt_sc(R,R_etoile) * alpha) / (pdt_sc(R_tmp,R_etoile) * omega);
        W = cl_vec(1.,R,beta,cl_vec(1.,W,-omega,AW));
        it++;
    }

    return X;
}


vector<double> erreurs(double (* solExa) (double,double, int),vector<double> uh,
    maillage TRG, int N, int M, double a, double b, int m){

    vector<double> tab_err(3);
    int G = (N+1)*(M+1);

    vector<double> uchap(G);

    for (int s = 0; s < G; s++){
        tuple<int,int> coord_noeud = invnumgb(N,M,s);
        double x = -a + (2*a)/N * get<0>(coord_noeud);
        double y = -b + (2*b)/M * get<1>(coord_noeud);

        uchap[s] = solExa(x,y,m);
    }


    vector<double> diff = cl_vec(1.,uchap,-1.,uh);

    tab_err[0] = normL2(diff,TRG,N,M,a,b) / normL2(uchap,TRG,N,M,a,b);
    tab_err[1] = normL2Grad(diff,TRG,N,M,a,b) / normL2Grad(uchap,TRG,N,M,a,b);
    tab_err[2] = max_abs(diff) / max_abs(uchap);

    return tab_err;
}
\end{codeblock}